% =============================================================================================
% DISCUSSION. Interpretation, practical guidance and limitations, moved out of Methods and Results
% (author decision 2026-10-11: separate Results and Discussion). Every number here already appears
% in Methods or Results and is quoted from data/exports/kappa_v2/paper_numbers.json.
% =============================================================================================

\section{Discussion}
\label{sec:discussion}

\subsection{What the offset and the constants mean}
\label{sec:interpretation}

Several effects can make published calculations run high, and they leave different fingerprints
(Table~\ref{tab:causes}). Real specimens scatter phonons at defects the calculation of a perfect
crystal lacks: interstitials and antisites, imaged directly in \ce{TiNiSn}-based
alloys~\cite{kim2017direct}, and grain boundaries~\cite{schrade2017the}; computed for \num{54}
half-Heusler alloys, alloy and grain-boundary scattering significantly reduce the intrinsic
conductivity~\cite{berland2019thermoelec}. Such scattering is nearly independent of temperature, so
it matters less as Umklapp scattering strengthens. Four-phonon scattering, which most calculations
omit, lowers the true conductivity, so omitting it biases three-phonon calculations
high~\cite{feng2017fourphonon}; a full treatment of quartic anharmonicity also shifts phonon
frequencies, however, and its net effect differs between structure types~\cite{xia2020highthroug}.
The exchange--correlation functional and force-constant cut-offs move a calculation either way. On
the measurement side, an electronic or bipolar term left in $\kappa_L$ by the Lorenz subtraction
raises the measured value~\cite{thesberg2017on,li2015synthesis}, and the bipolar part grows with
temperature as more carriers are excited across the gap.

\begin{table}[!b]
  \centering
  \caption{Possible causes of the offset and the signature each would leave. ``Ratio'' is
  calculated/measured.}
  \label{tab:causes}
  \footnotesize
  \setlength{\tabcolsep}{3pt}
  \begin{tabular}{>{\raggedright\arraybackslash}p{0.42\columnwidth}>{\raggedright\arraybackslash}p{0.18\columnwidth}>{\raggedright\arraybackslash}p{0.30\columnwidth}}
    \toprule
    Cause & Ratio & As $T$ rises \\
    \midrule
    defects, grain boundaries, pores in the specimen & above 1 & shrinks towards 1 \\
    four-phonon scattering left out of the calculation & usually above 1 & grows \\
    functional, cut-offs, convergence & either & weakly dependent \\
    electronic or bipolar term left in measured $\kappa_L$ & lowers it & effect grows \\
    \bottomrule
  \end{tabular}
\end{table}

Read against Table~\ref{tab:causes}, the three observations of Section~\ref{sec:offset} narrow the
possibilities without settling them. The offset shrinks with temperature, which fits specimen
scattering or a residual electronic term better than missing four-phonon scattering, whose effect
would grow. It exceeds the spread between independent calculations, so differences between codes
and settings are not enough to explain it. And at \SI{300}{\kelvin} it shows no trend with the
calculated conductivity: a single boundary mean free path common to every specimen would suppress
the most conductive compounds, whose intrinsic mean free paths are longest, more strongly. That
argues against one common grain size as the whole explanation; defect populations that differ from
specimen to specimen, or errors in the calculations themselves, remain possible. For scale, if the
whole \SI{300}{\kelvin} offset were extrinsic scattering added by Matthiessen's rule,
$\kappa_L^{\mathrm{exp}}/\kappa_L^{\mathrm{BTE}} = 1/(1+\tau_{\mathrm{int}}/\tau_{\mathrm{ext}})$ in a
grey model, the bin median of \num{1.84} would require an extrinsic scattering rate \num{0.84} times
the intrinsic one.

Phenomenological temperature dependences have been fitted to, or tested against, measured
half-Heusler curves before~\cite{geng2014lattice,petersen2015critical}; here we fit instead the \emph{ratio} of
measurement to calculation. The functional form is therefore not new; what this work adds is the
corpus, the class-wide magnitude and its held-out validation. Because nearly every pair carries the
assumed $T^{-1}$ calculated dependence (Section~\ref{sec:transfer}; the one multi-temperature
calculation has a log--log slope of \num{-1.02}), $p$ is in effect the measured slope relative to
it: $p=\num{0.70}$ implies a corrected $\kappa_L \propto T^{-0.30}$.

The correction is empirical, not a boundary-scattering term. In a relaxation-time description of
the conductivity~\cite{callaway1959model}, any temperature-independent extrinsic scattering
(boundaries, point defects or static disorder) added to three-phonon Umklapp scattering caps the
logarithmic slope of the suppression at $1-c$ (\num{0.58} at $c=\num{0.42}$), for a grey or a
spectral model; this bound is our own derivation (Section~S1 of the Supplementary material). The
fitted $p$ exceeds that cap, although its bootstrap range straddles it, and four-phonon scattering
would lower the slope. A measured lattice conductivity above the three-phonon value of
the perfect crystal, as some high-temperature curves show, more likely retains an electronic or
bipolar term that the single-parabolic-band Lorenz number does not remove, or high-temperature
measurement error: the Lorenz number of a multiband or bipolar semiconductor can depart from the
single-band value by a factor of two or more~\cite{thesberg2017on}, and a bipolar contribution to the
thermal conductivity has been reported in a half Heusler of this domain~\cite{li2015synthesis}. The rise of the measured-to-calculated ratio with temperature, and the value of
$p$, may therefore partly reflect the separation of lattice and electronic conduction rather than
phonon scattering. We attach no mechanism to the parameters.

A family constant describes the specimens behind it, and it may also absorb measurement effects. The
\ce{Bi}--\ce{Pd} constant is close to one, and its rare-earth members are the narrow-gap compounds for
which the electronic subtraction is least secure (Section~\ref{sec:domain}); part of it may reflect an
electronic term left in the measured values rather than better crystals. Several families rest on a
single kind of
specimen (the publications behind each family are counted in Section~S5 of the Supplementary
material). All three \ce{Sn}--\ce{Pt} references come from one paper, as do three of the four in
\ce{Sb}--\ce{Pt}, four of the seven in \ce{Ni}--\ce{Sb} and two of the three in each of
\ce{Bi}--\ce{Pd} and \ce{Sb}--\ce{Pd}; there the held-out test partly measures one paper's internal
consistency rather than transfer between laboratories. Two \ce{Sb}--\ce{Pd} references are
annealed arc-melted ingots whose authors attribute the near-linear low-temperature conductivity to
additional scattering at grain boundaries, lattice defects and
disorder~\cite{mukhopadhyay2018multifunct}. The two \ce{Co}--\ce{Sn} members were measured on very
different specimens, \ce{NbCoSn} on a near-fully dense pellet~\cite{thiem2025thermoelec} and
\ce{TaCoSn} on a heavily ball-milled nanocrystalline, cobalt-rich sample that its authors suggest
contains cobalt interstitials~\cite{li2019ntype}, and a constant fitted on one cannot describe the other. An
anchored family inherits its single member's specimen. \ce{ZrNiPb} was measured only on a
ball-milled sample~\cite{mao2017thermoelec}, so values issued through \ce{Ni}--\ce{Pb} describe a
milled sample. The \ce{ZrCoBi} reference is a \SI{98.1}{\percent} dense
pellet~\cite{zhao2018synthesis}, whereas the only other laboratory to measure it reports a value
\num{1.7} times lower near \SI{300}{\kelvin} on a far more heavily milled
specimen~\cite{zhu2018discovery}. The \ce{Co}--\ce{Bi} predictions therefore carry a wider uncertainty
than their conformal band.

\subsection{Which route to use}
\label{sec:routes}

The results suggest a simple order of preference for estimating the measured $\kappa_L$ of an
unmeasured half Heusler. Where a full transport calculation has been published and the compound's
bonding family has been validated, the family constant is the most accurate option tested
(\SI{21.7}{\percent} median error held out, against \SI{32.3}{\percent} for the shared constant on
the same compounds). Where a calculation exists but the family has no measured member, only the
shared constant applies, with the \SI{33.0}{\percent} median error measured on families withheld
from the fit. Where no calculation exists, the machine-learned model supplies one, and the value
carries two separate errors: \SI{21.9}{\percent} for the model step (reference seed) and \SI{33.0}{\percent} for the
shared constant.

This order has two limits. A family that has already been measured well can be described
almost as accurately by a power law fitted to its own members (\SI{25.4}{\percent} against
\SI{31.6}{\percent} for the model route, not significantly different), so the method adds most
where measurements are sparse: in ordering compounds within a family and in families never
measured. And no route predicts a temperature dependence; every value is a value at a stated
temperature.

\subsection{Limitations}
\label{sec:limitations}

Temperature dependence is not predicted. One exponent serves every compound and it is poorly
determined (bootstrap range \numrange{0.45}{1.15} against \numrange{0.35}{0.55} for $c$;
Section~\ref{sec:transfer}), and the training labels sit mostly at \SI{300}{\kelvin}. Compound by
compound, the predicted log--log slopes do not track the measured ones: on \num{33} compounds with a
physical measured slope, the measured slopes span \numrange{-4.78}{-0.02} and the predicted ones
\numrange{-0.27}{0.64}, with Spearman $\rho = \num{0.22}$ ($p = \num{0.22}$). For some compounds the
route therefore predicts a conductivity that rises with temperature, whereas all \num{33} measured
slopes are negative; every value in this paper is accordingly quoted at one stated temperature. At
high
temperature the measured lattice values may also retain part of the electronic contribution, which
then enters the fitted exponent (Section~\ref{sec:interpretation}).

Each compound is scored against one reference laboratory. Laboratories measuring the same in-domain
compound disagree by a median factor of \num{1.33} over all pairs of measurements (\num{1.36} taking
each compound's median), so a median error of about a third is comparable in size to the
reproducibility of the field, although the two are different measures. Per-family errors may fall
below that level partly because each compound is scored against one chosen reference.

The compounds are not independent. They share chemistry clusters, bonding families and, in places, a
laboratory, so the effective sample is smaller than \num{39}; the cluster-level tests address the
first, not the shared laboratories.

A family constant is interpolation within the family, and in several families the family-mates are
one paper's specimens. A family nobody has measured receives only the shared constant, and the
capped shared correction cannot raise a calculation that already lies below its measurement.

The issued values rest on few calculations. Seven of the ten rest on a single published calculation
at the quoted temperature, and most on one study. A systematic error there would pass undiluted into
those values; a transfer function calibrates the scale of a calculation, not its correctness.

Every half-Heusler training label is a full calculation, which leaves \num{284} compounds to learn
from;
where a bond descriptor cannot be computed it is left missing, which the model handles natively.
Structure-only values carry two errors and assume a structure: neither compound has, as far as we
could find, been synthesised, so each value is conditional on the cubic phase forming.

Several alternatives were tested and not adopted, because they did not improve held-out accuracy
consistently: a correction without the temperature term (Section~\ref{sec:robust}), a capped family
form, a two-stage (hurdle) model, a magnitude-dependent calibration term, and training on
experimental or doped-composition values.
